\hme{move to global prelims}

We use $[x,  y]$ to denote the set $\{x, x+1,\dots,y\}$, where $x, y \in \ZZ$ and $x \leq y$.
We use $[x]$ as shorthand to denote $[1, x]$.
When $x > y$ we allow $[x, y]$ to denote $\{x, x-1,\dots, y\}$.
For an integer $x$ we use $\xthdigit{x}{i}_b$ to denote the $1$-indexed $i$'th digit of $x$ in base $b$, e.g. for the base-$10$ value $x = 357$, $\xthdigit{x}{1}_{10} = 3$.
We may drop the subscript when the base is clear from context.
We denote concatenation between the strings $x$ and $y$ with $x ||y$, and the empty string as $\varepsilon$.
For any algorithm $\adv$, we use $y := \adv(x; r)$ to denote that the algorithm runs on input $x$ and random tape $r$.
In most cases, we do not denote the random tape explicitly and simply write $y \leftarrow \adv(x)$.
We may write $y \xleftarrow{r} \adv(x)$ to indicate sampling $r$ as appropriate and recording it for later use.
We use $\secpar$ as the computational security parameter, and $\negl(\secpar)$ to denote negligible functions.
For two probability ensembles $\{A_i\}_i$ and $\{B_i\}_i$, we use $\{A_i\}_i \pindist \{B_i\}_i$ to denote that the ensembles are identical, $\{A_i\}_i \sindist \{B_i\}_i$ to denote that the ensembles are statistically close, and $\{A_i\}_i \cindist \{B_i\}_i$ to denote that the ensembles are computationally indistinguishable. We use $\prob{E : A}$ to denote the probability of an event $E$ in an experiment defined by executing $A$.
For a finite set $\mathcal{X}$, we will use $\mathcal{U}_{\mathcal{X}}$ to denote the uniform distribution over $\mathcal{X}$.
For a distribution $\dist$ we will use $x \leftarrow \dist$ to indicate sampling a value according to the distribution.
We will use $\fthr$ to denote the following function:
\fthrleak

Where $f'_b(v, t)$ is defined as follows.
\procedureblock{$f'_b(\uval, \thr)$}{
	x = \varepsilon \\
	\pcfor i \in [\lceil \log_\base(\thr)\rceil] \\
	\t \pcif \uval[i]_b \geq \thr[i]_b \pcthen x = x || \uval[i]_b
	\: \pcelse \textbf{break}\\
	\pcreturn x
}


